
\author{W. E. B. Du Bois}
\authdate{1868--1963}
\textdate{1897/1903}
\chapter{Of Our Spiritual Strivings}
\source{dubois1903.2}

% First published in The Atlantic Monthly in August 1897, under the
% title ``Strivings of the Negro People.''

\settowidth\versewidth{\small\vin\vin The voice of my heart in my side
or the voice of the sea,}
\page{1}\begin{verse}[\versewidth]
O water, voice of my heart, crying in the sand,\\
\vin All night long crying with a mournful cry,\\
As I lie and listen, and cannot understand\\
\vin\vin The voice of my heart in my side or the voice of the sea,\\
\vin O water, crying for rest, is it I, is it I?\\
\vin\vin All night long the water is crying to me.

Unresting water, there shall never be rest\\
\vin Till the last moon droop and the last tide fail,\\
And the fire of the end begin to burn in the west;\\
\vin\vin And the heart shall be weary and wonder and cry like the sea,\\
\vin All life long crying without avail,\\
\vin\vin As the water all night long is crying to me.
\end{verse}
\attrib{Arthur Symons.}

\begin{center}
\includegraphics[width=\maxwidth]{texts/music.eps}
\end{center}

\noindent Between me and the other world there is ever an unasked
question: unasked by some through feelings of delicacy; by others
through the difficulty of rightly framing it. All, nevertheless,
flutter round it. They approach me in a half-hesitant sort of way, eye
me curiously or compassionately, and then, instead of saying directly,
How does it feel to be a problem? they say, I know an excellent
\page{2} colored man in my town; or, I fought at Mechanicsville; or,
Do not these Southern outrages make your blood boil? At these I smile,
or am interested, or reduce the boiling to a simmer, as the occasion
may require. To the real question, How does it feel to be a problem? I
answer seldom a word.

And yet, being a problem is a strange experience,---pe\-cu\-liar even
for one who has never been anything else, save perhaps in babyhood and
in Europe. It is in the early days of rollicking boyhood that the
revelation first bursts upon one, all in a day, as it were. I
remember well when the shadow swept across me. I was a little thing,
away up in the hills of New England, where the dark Housatonic winds
between Hoosac and Taghkanic to the sea. In a wee wooden schoolhouse,
something put it into the boys' and girls' heads to buy gorgeous
visiting-cards---ten cents a package---and exchange. The exchange was
merry, till one girl, a tall newcomer, refused my card,---re\-fused it
peremptorily, with a glance. Then it dawned upon me with a certain
suddenness that I was different from the others; or like, mayhap, in
heart and life and longing, but shut out from their world by a vast
veil. I had thereafter no desire to tear down that veil, to creep
through; I held all beyond it in common contempt, and lived above it
in a region of blue sky and great wandering shadows. That sky was
bluest when I could beat my mates at examination-time, or beat them at
a foot-race, or even beat their stringy heads. Alas, with the years
all this fine contempt began to fade; for the worlds I longed for, and
all their dazzling opportunities, were \page{3} theirs, not mine. But
they should not keep these prizes, I said; some, all, I would wrest
from them. Just how I would do it I could never decide: by reading
law, by healing the sick, by telling the wonderful tales that swam
in my head,---some way. With other black boys the strife was not so
fiercely sunny: their youth shrunk into tasteless sycophancy, or into
silent hatred of the pale world about them and mocking distrust of
everything white; or wasted itself in a bitter cry, Why did God make
me an outcast and a stranger in mine own house? The shades of the
prison-house closed round about us all: walls strait and stubborn to
the whitest, but relentlessly narrow, tall, and unscalable to sons
of night who must plod darkly on in resignation, or beat unavailing
palms against the stone, or steadily, half hopelessly, watch the
streak of blue above.

After the Egyptian and Indian, the Greek and Roman, the Teuton and
Mongolian, the Negro is a sort of seventh son, born with a veil, and
gifted with second-sight in this American world,---a world which
yields him no true self-consciousness, but only lets him see himself
through the revelation of the other world. It is a peculiar sensation,
this double-consciousness, this sense of always looking at one's self
through the eyes of others, of measuring one's soul by the tape of a
world that looks on in amused contempt and pity. One ever feels his
twoness,---an American, a Negro; two souls, two thoughts, two
unreconciled strivings; two warring ideals in one dark body, whose
dogged strength alone keeps it from being torn asunder.

\page{4}The history of the American Negro is the history of this
strife,---this longing to attain self-conscious manhood, to merge his
double self into a better and truer self. In this merging he wishes
neither of the older selves to be lost. He would not Africanize
America, for America has too much to teach the world and Africa. He
would not bleach his Negro soul in a flood of white Americanism, for
he knows that Negro blood has a message for the world. He simply
wishes to make it possible for a man to be both a Negro and an
American, without being cursed and spit upon by his fellows, without
having the doors of Opportunity closed roughly in his face.

This, then, is the end of his striving: to be a co-worker in the
kingdom of culture, to escape both death and isolation, to husband and
use his best powers and his latent genius. These powers of body and
mind have in the past been strangely wasted, dispersed, or forgotten.
The shadow of a mighty Negro past flits through the tale of Ethiopia
the Shadowy and of Egypt the Sphinx. Through history, the powers of
single black men flash here and there like falling stars, and die
sometimes before the world has rightly gauged their brightness. Here
in America, in the few days since Emancipation, the black man's
turning hither and thither in hesitant and doubtful striving has often
made his very strength to lose effectiveness, to seem like absence of
power, like weakness. And yet it is not weakness,---it is the
contradiction of double aims. The double-aimed struggle of the black
artisan---on the one hand to escape white contempt for a nation of
mere hewers of wood and drawers \page{5} of water, and on the other
hand to plough and nail and dig for a poverty-stricken horde---could
only result in making him a poor craftsman, for he had but half a
heart in either cause. By the poverty and ignorance of his people, the
Negro minister or doctor was tempted toward quackery and demagogy; and
by the criticism of the other world, toward ideals that made him
ashamed of his lowly tasks. The would-be black \textit{savant} was
confronted by the paradox that the knowledge his people needed was a
twice-told tale to his white neighbors, while the knowledge which
would teach the white world was Greek to his own flesh and blood. The
innate love of harmony and beauty that set the ruder souls of his
people a-dancing and a-singing raised but confusion and doubt in the
soul of the black artist; for the beauty revealed to him was the
soul-beauty of a race which his larger audience despised, and he could
not articulate the message of another people. This waste of double
aims, this seeking to satisfy two unreconciled ideals, has wrought sad
havoc with the courage and faith and deeds of ten thousand thousand
people,---has sent them often wooing false gods and invoking false
means of salvation, and at times has even seemed about to make them
ashamed of themselves.

Away back in the days of bondage they thought to see in one divine
event the end of all doubt and disappointment; few men ever worshipped
Freedom with half such unquestioning faith as did the American Negro
for two centuries. To him, so far as he thought and dreamed, slavery
was indeed the sum of all villainies, the cause of all sorrow, the
root of all \page{6} prejudice; Emancipation was the key to a promised
land of sweeter beauty than ever stretched before the eyes of wearied
Israelites. In song and exhortation swelled one refrain---Liberty; in
his tears and curses the God he implored had Freedom in his right
hand. At last it came,---suddenly, fearfully, like a dream. With one
wild carnival of blood and passion came the message in his own
plaintive cadences:---

\begin{verse}
\poke{``}Shout, O children!\\
Shout, you're free!\\
For God has bought your liberty!''
\end{verse}

Years have passed away since then,---ten, twenty, forty; forty years
of national life, forty years of renewal and development, and yet the
swarthy spectre sits in its accustomed seat at the Nation's feast. In
vain do we cry to this our vastest social problem:---

\begin{verse}
\poke{``}Take any shape but that, and my firm nerves\\
Shall never tremble!''
\end{verse}

The Nation has not yet found peace from its sins; the freedman has not
yet found in freedom his promised land. Whatever of good may have come
in these years of change, the shadow of a deep disappointment rests
upon the Negro people,---a disappointment all the more bitter because
the unattained ideal was unbounded save by the simple ignorance of a
lowly people.

The first decade was merely a prolongation of the vain search for
freedom, the boon that seemed ever barely to elude their grasp,---like
a tantalizing will-o'-the-wisp, maddening and misleading the
head-\page{7}less host. The holocaust of war, the terrors of the
Ku-Klux Klan, the lies of carpet-baggers, the disorganization of
industry, and the contradictory advice of friends and foes, left the
bewildered serf with no new watchword beyond the old cry for freedom.
As the time flew, however, he began to grasp a new idea. The ideal of
liberty demanded for its attainment powerful means, and these the
Fifteenth Amendment gave him. The ballot, which before he had looked
upon as a visible sign of freedom, he now regarded as the chief means
of gaining and perfecting the liberty with which war had partially
endowed him. And why not? Had not votes made war and emancipated
millions? Had not votes enfranchised the freedmen? Was anything
impossible to a power that had done all this? A million black men
started with renewed zeal to vote themselves into the kingdom. So the
decade flew away, the revolution of 1876 came, and left the half-free
serf weary, wondering, but still inspired. Slowly but steadily, in the
following years, a new vision began gradually to replace the dream of
political power,---a powerful movement, the rise of another ideal to
guide the unguided, another pillar of fire by night after a clouded
day. It was the ideal of ``book-learning''; the curiosity, born of
compulsory ignorance, to know and test the power of the cabalistic
letters of the white man, the longing to know. Here at last seemed to
have been discovered the mountain path to Canaan; longer than the
highway of Emancipation and law, steep and rugged, but straight,
leading to heights high enough to overlook life.

\page{8}Up the new path the advance guard toiled, slowly, heavily,
doggedly; only those who have watched and guided the faltering feet,
the misty minds, the dull understandings, of the dark pupils of these
schools know how faithfully, how piteously, this people strove to
learn. It was weary work. The cold statistician wrote down the inches
of progress here and there, noted also where here and there a foot had
slipped or some one had fallen. To the tired climbers, the horizon was
ever dark, the mists were often cold, the Canaan was always dim and
far away. If, however, the vistas disclosed as yet no goal, no
resting-place, little but flattery and criticism, the journey at least
gave leisure for reflection and self-examination; it changed the child
of Emancipation to the youth with dawning self-consciousness,
self-realization, self-respect. In those sombre forests of his
striving his own soul rose before him, and he saw
him\-self,---dark\-ly as through a veil; and yet he saw in himself
some faint revelation of his power, of his mission. He began to have a
dim feeling that, to attain his place in the world, he must be
himself, and not another. For the first time he sought to analyze the
burden he bore upon his back, that dead-weight of social degradation
partially masked behind a half-named Negro problem. He felt his
poverty; without a cent, without a home, without land, tools, or
savings, he had entered into competition with rich, landed, skilled
neighbors. To be a poor man is hard, but to be a poor race in a land
of dollars is the very bottom of hardships. He felt the weight of his
ignorance,---not simply of letters, but of life, of \page{9} business,
of the humanities; the accumulated sloth and shirking and awkwardness
of decades and centuries shackled his hands and feet. Nor was his
burden all poverty and ignorance. The red stain of bastardy, which two
centuries of systematic legal defilement of Negro women had stamped
upon his race, meant not only the loss of ancient African chastity,
but also the hereditary weight of a mass of corruption from white
adulterers, threatening almost the obliteration of the Negro home.

A people thus handicapped ought not to be asked to race with the
world, but rather allowed to give all its time and thought to its own
social problems. But alas! while sociologists gleefully count his
bastards and his prostitutes, the very soul of the toiling, sweating
black man is darkened by the shadow of a vast despair. Men call the
shadow prejudice, and learnedly explain it as the natural defence of
culture against barbarism, learning against ignorance, purity against
crime, the ``higher'' against the ``lower'' races. To which the Negro
cries Amen! and swears that to so much of this strange prejudice as is
founded on just homage to civilization, culture, righteousness, and
progress, he humbly bows and meekly does obeisance. But before that
nameless prejudice that leaps beyond all this he stands helpless,
dismayed, and well-nigh speechless; before that personal disrespect
and mockery, the ridicule and systematic humiliation, the distortion
of fact and wanton license of fancy, the cynical ignoring of the
better and the boisterous welcoming of the worse, the all-pervading
desire to inculcate disdain for everything \page{10} black, from
Toussaint to the devil,---be\-fore this there rises a sickening
despair that would disarm and discourage any nation save that black
host to whom ``discouragement'' is an unwritten word.

But the facing of so vast a prejudice could not but bring the
inevitable self-questioning, self-disparagement, and lowering of
ideals which ever accompany repression and breed in an atmosphere of
contempt and hate. Whisperings and portents came borne upon the four
winds: Lo! we are diseased and dying, cried the dark hosts; we cannot
write, our voting is vain; what need of education, since we must
always cook and serve? And the Nation echoed and enforced this
self-criticism, saying: Be content to be servants, and nothing more;
what need of higher culture for half-men? Away with the black man's
ballot, by force or fraud,---and behold the suicide of a race!
Nevertheless, out of the evil came something of good,---the more
careful adjustment of education to real life, the clearer perception
of the Negroes' social responsibilities, and the sobering realization
of the meaning of progress.

So dawned the time of \textit{Sturm und Drang}: storm and stress
to-day rocks our little boat on the mad waters of the world-sea; there
is within and without the sound of conflict, the burning of body and
rending of soul; inspiration strives with doubt, and faith with vain
questionings. The bright ideals of the past,---physical freedom,
political power, the training of brains and the training of
hands,---all these in turn have waxed and waned, until even the last
grows dim and overcast. Are they all wrong,---all false? No, not that,
but \page{11} each alone was over-simple and incomplete,---the dreams
of a credulous race-childhood, or the fond imaginings of the other
world which does not know and does not want to know our power. To be
really true, all these ideals must be melted and welded into one. The
training of the schools we need to-day more than ever,---the training
of deft hands, quick eyes and ears, and above all the broader, deeper,
higher culture of gifted minds and pure hearts. The power of the
ballot we need in sheer self-defence,---else what shall save us from a
second slavery? Freedom, too, the long-sought, we still seek,---the
freedom of life and limb, the freedom to work and think, the freedom
to love and aspire. Work, culture, liberty,---all these we need, not
singly but together, not successively but together, each growing and
aiding each, and all striving toward that vaster ideal that swims
before the Negro people, the ideal of human brotherhood, gained
through the unifying ideal of Race; the ideal of fostering and
developing the traits and talents of the Negro, not in opposition to
or contempt for other races, but rather in large conformity to the
greater ideals of the American Republic, in order that some day on
American soil two world-races may give each to each those
characteristics both so sadly lack. We the darker ones come even now
not altogether empty-handed: there are to-day no truer exponents of
the pure human spirit of the Declaration of Independence than the
American Negroes; there is no true American music but the wild sweet
melodies of the Negro slave; the American fairy tales and folk-lore
are Indian and African; and, all in all, we black men seem the sole
\page{12} oasis of simple faith and reverence in a dusty desert of
dollars and smartness. Will America be poorer if she replace her
brutal dyspeptic blundering with light-hearted but determined Negro
humility? or her coarse and cruel wit with loving jovial good-humor?
or her vulgar music with the soul of the Sorrow Songs?

Merely a concrete test of the underlying principles of the great
republic is the Negro Problem, and the spiritual striving of the
freedmen's sons is the travail of souls whose burden is almost beyond
the measure of their strength, but who bear it in the name of an
historic race, in the name of this the land of their fathers' fathers,
and in the name of human opportunity.

\vspace{1\baselineskip}

And now what I have briefly sketched in large outline let me on coming
pages tell again in many ways, with loving emphasis and deeper detail,
that men may listen to the striving in the souls of black folk.

